\chapter{INTRODUCTION}

\section{Background}

Cardiovascular diseases (CVDs) and cardiac rhythm disorders remain the leading cause of mortality worldwide, accounting for an estimated 17.9 million deaths annually \cite{who2025cvd}. Early detection of cardiac rhythm abnormalities---including ectopic contractions, irregular atrial rhythms, and resting rate disorders---is critical because irregular beat-to-beat timing strongly correlates with stroke, heart failure, and sudden cardiac events. In clinical practice, 12-lead electrocardiography (ECG) and ambulatory Holter monitors represent the diagnostic gold standard; however, both rely on specialized analog front-ends, conductive gels, and adhesive chest electrodes that create cost, discomfort, and practical barriers to frequent out-of-clinic screening. Consumer smartwatches partially lower this barrier, yet they still require purchasing dedicated \$150--\$400+ wearable hardware and rely on wrist-mounted optical photoplethysmography (PPG). Public validation benchmarks demonstrate that peripheral wrist PPG suffers from skin-pigmentation attenuation, poor peripheral perfusion, motion-induced error spikes, and vascular Pulse Transit Time (PTT) jitter \cite{bent2020investigating, nelson2019accuracy}. Consequently, there is a compelling engineering need for a zero-external-hardware screening modality built directly on consumer smartphones already owned by the general public.

Sternal seismocardiography (SCG) offers a direct, non-invasive alternative that captures central cardiac mechanical activity without electrical leads or external peripherals. Specifically, SCG records the micro-vibrations transmitted to the anterior chest wall by ventricular ejection recoil and aortic valve opening (AO) \cite{inan2015ballistocardiography, tadi2016realtime}. By placing an unmodified smartphone flat over the lower third of the sternum during resting supine posture, the device's built-in three-axis micro-electromechanical system (MEMS) accelerometer can measure these dorso-ventral ($Z$-axis) vibrations directly. This creates a clear engineering distinction and selling point relative to smartwatches: smartphone sternal SCG operates at \textbf{\$0 incremental hardware cost} and directly observes \textbf{central mechanical rhythm} rather than peripheral optical blood volume. Moreover, rather than competing with wearables during exercise, smartphone sternal SCG enables an optional \textbf{Complementary Dual-Mode Pathway} in which smartwatches handle continuous 24/7 ambulatory alerts \cite{perez2019large, lubitz2022detection} while a smartphone provides zero-cost, on-demand central mechanical confirmation at rest.

Translating raw smartphone MEMS acceleration into clinically meaningful cardiac metrics presents substantial signal-processing challenges. Consumer accelerometers suffer from higher thermal noise floors, DC gravity offsets, low-frequency respiratory baseline drift, high-frequency muscular tremor, posture misalignment, and mobile operating system timestamp jitter. Robust digital signal processing (DSP) is therefore essential: zero-phase bandpass filtering must isolate the cardiac mechanical passband without distorting peak timestamps, while automated quality gating must reject stationary inanimate surfaces and motion-corrupted recordings before peak detection occurs. From verified AO systolic peaks, Peak-to-Peak Interval (PPI) tachograms are constructed to extract standardized Pulse Rate Variability (PRV) features---spanning time-domain metrics (Mean HR, SDNN, RMSSD, pNN50) and frequency-domain spectral powers (LF, HF, LF/HF) defined by the Task Force of ESC and NASPE \cite{taskforce1996heart}. As established by Sch{\"a}fer and Vagedes \cite{schafer2013accurate}, PRV derived from mechanical surrogates closely approximates ECG-derived heart rate variability under stationary resting conditions when motion artifacts are rigorously controlled.

Prior research establishes strong individual foundations for this work. Inan et al. \cite{inan2015ballistocardiography} and Tadi et al. \cite{tadi2016realtime} demonstrated that chest-mounted and smartphone sternal MEMS accelerometers can reliably localize AO peaks and track resting heart rate. Tsamis et al. \cite{tsamis2017soft} analyzed soft real-time execution budgets and buffer management for mobile cardiac DSP, while Aliev et al. \cite{aliev2012monitoring} proved that noise-variance gating substantially reduces false clinical alarms in mobile monitoring. In wearable validation, Perez et al. \cite{perez2019large}, Lubitz et al. \cite{lubitz2022detection}, Bent et al. \cite{bent2020investigating}, and Nelson and Allen \cite{nelson2019accuracy} established public smartwatch PPG performance benchmarks for resting HR and atrial fibrillation screening. In dynamical modeling, van der Pol and van der Mark \cite{vanderpol1928} and Quiroz-Ju{\'a}rez et al. \cite{quiroz2018cardiac} formulated coupled nonlinear cardiac oscillator models capable of reproducing physiological and arrhythmic rhythm dynamics. Despite these advances, three practical engineering gaps limit current mobile cardiac screening tools:
\begin{enumerate}
	\item \textbf{Hardware Cost and Peripheral Optical Limitations:} Existing consumer solutions depend on external ECG peripherals or costly wrist-worn PPG smartwatches whose peripheral optical sensing introduces skin-tone, perfusion, and PTT distortion \cite{bent2020investigating}, whereas existing smartphone apps rarely benchmark their central mechanical accuracy against public wearable data.
	\item \textbf{Lack of Active Mechanical Quality Gating and Guided Usability:} Consumer MEMS sensors easily produce false heartbeat detections when resting on vibrating inanimate surfaces or when the user moves or tilts the device. Reliable zero-hardware sternal SCG requires deterministic physiological energy boundaries, SNR verification, and real-time posture and stillness guidance that is measurably fast and intuitive for non-technical users.
	\item \textbf{Opaque Risk Scores, Leakage-Prone Validation, and Pathological SCG Scarcity:} Many mobile health tools output black-box scores validated on record-mixed splits rather than strict patient-stratified benchmarks \cite{dechazal2004automatic}. Furthermore, acquiring real clinical arrhythmia recordings directly on smartphones is constrained by hospital ethics and patient recruitment, creating a need for a controlled synthetic SCG simulation bridge to verify end-to-end screening sensitivity.
\end{enumerate}

To address these gaps, this senior capstone project proposes an integrated, zero-hardware smartphone sternal SCG platform focused on a single physiological sensing domain---resting beat-to-beat PPI/PRV timing and heart rate---scoped specifically to a \textbf{Rhythm-Disorder Cluster under Cardiac Arrhythmia}: (1) \textbf{Ectopic Beats} (Premature Ventricular Contractions [PVC, AAMI $V$] and Premature Atrial/Supraventricular Contractions [PAC, AAMI $S$]), (2) \textbf{Irregular Rhythm / Atrial Fibrillation (AFib)}, and (3) \textbf{Rate Disorders} (resting Sinus Bradycardia $<60\text{ bpm}$ and Sinus Tachycardia $>100\text{ bpm}$). Because all three sub-conditions manifest directly in beat-to-beat interval timing without requiring 12-lead ECG morphology, they share one unified feature space. The methodology follows a progressive three-stage validation workflow: (1) acquiring resting lower-sternal SCG signals via the internal MEMS accelerometer with zero-phase filtering, physiological energy gating, guided usability feedback, and lightweight moving-RMS AO peak detection; (2) verifying inter-subject physiological differentiation across distinct individuals and public benchmarks; and (3) executing \textbf{Two-Stage Hybrid Screening}---combining deterministic physiological HR rules for Rate Disorders with explainable edge machine learning classifiers (\texttt{LR}, \texttt{RF}, \texttt{GBDT/XGBoost}) benchmarked against a deep spatio-temporal comparator (\texttt{Hybrid 1D-CNN + BiLSTM + Bahdanau Attention}) under strict patient-stratified clinical datasets, supported by a coupled Van der Pol oscillator serving as a controlled arrhythmia simulation tool.

\section{Objective}
Every engineering objective in this capstone is formulated with explicit, verifiable quantitative targets spanning \textbf{Accuracy}, \textbf{Efficiency}, and \textbf{User-Friendliness}:
\begin{enumerate}
	\item \textbf{To develop} a zero-external-hardware mobile sternal SCG acquisition pipeline using a consumer smartphone's internal 3-axis MEMS accelerometer ($100\text{--}200\text{ Hz}$) with real-time posture ($\pm 10^\circ$) and stillness guidance, targeting \textbf{User-Friendliness} metrics of guided setup completion in $< 60\text{ s}$, first-attempt valid-signal capture rate $\ge 90\%$, and a System Usability Scale score $\text{SUS} \ge 75$ \cite{brooke1996sus}.
	\item \textbf{To develop} a deterministic zero-phase digital signal processing (DSP) engine (4th-order Butterworth $0.5\text{--}40\text{ Hz}$ bandpass, physiological energy gate $0.005 \le \sigma_{raw} \le 0.50\text{ m/s}^2$, $\text{SNR} \ge 12\text{ dB}$, and $200\text{ ms}$ moving-RMS AO peak detector) that achieves \textbf{Efficiency} of $< 500\text{ ms}$ end-to-end processing latency per window, $100\%$ rejection of stationary tabletop and active motion artifacts, and \textbf{Accuracy} of reference heart-rate $\text{MAE} < 3.0\text{ bpm}$ and $\text{RMSE} < 4.5\text{ bpm}$.
	\item \textbf{To verify} inter-subject physiological differentiation by comparing extracted PPI tachograms and PRV feature distributions across at least two distinct healthy subjects (empirical smartphone SCG) and patient-stratified public clinical records, confirming statistically repeatable individual separability rather than generic sensor noise.
	\item \textbf{To build and benchmark} a \textbf{Two-Stage Hybrid Screening} engine for the \textbf{Rhythm-Disorder Cluster}---flagging Rate Disorders (Bradycardia $< 60\text{ bpm}$, Tachycardia $> 100\text{ bpm}$) via deterministic rules and classifying Rhythm Arrhythmias (Normal Sinus Rhythm vs. Ectopic [PVC/PAC] vs. Irregular/AFib, plus a Binary Abnormal flag) via explainable edge models (Logistic Regression, Random Forest, GBDT/XGBoost; model footprint $< 2\text{ MB}$) benchmarked against a \texttt{Hybrid 1D-CNN + BiLSTM + Bahdanau Attention} reference on patient-stratified clinical datasets (MIT-BIH DS1/DS2 \cite{goldberger2000physiobank, dechazal2004automatic}, INCART, PTB-XL, CPSC 2018, Chapman-Shaoxing)---targeting overall $\text{Accuracy} \ge 85\%$, $\text{Sensitivity} \ge 88\%$, $\text{Specificity} \ge 85\%$, and $\text{F1-score} \ge 0.86$.
	\item \textbf{To simulate} controlled arrhythmic SCG test sequences using a coupled Van der Pol oscillator (RK45) combined with empirical sternal AO beat templates to verify end-to-end pipeline sensitivity under realistic MEMS noise, and \textbf{to benchmark} smartphone sternal SCG performance quantitatively against published public smartwatch PPG datasets \cite{bent2020investigating, nelson2019accuracy, perez2019large, lubitz2022detection}.
\end{enumerate}

\section{Scope}
\begin{enumerate}
	\item \textbf{What you do or do not consider in this project:}
	\begin{itemize}
		\item \textit{Inclusions:} Non-invasive sternal SCG acquisition using standard consumer smartphone 3-axis MEMS accelerometers ($\pm 2g$, $100\text{--}200\text{ Hz}$) during resting supine postures; optimal lower-sternum placement validation; 4th-order zero-phase Butterworth bandpass filtering ($0.5\text{--}40\text{ Hz}$); physiological energy gating ($\sigma_{min} = 0.005\text{ m/s}^2$, $\sigma_{max} = 0.50\text{ m/s}^2$) and SNR verification ($\ge 12\text{ dB}$); $200\text{ ms}$ moving-RMS envelope and adaptive threshold AO peak detection; time-domain PRV extraction (Mean HR, SDNN, RMSSD, pNN50) and frequency-domain spectral analysis (LF, HF, LF/HF) via $4.0\text{ Hz}$ cubic-spline resampling and Welch FFT; 2-person inter-subject physiological differentiation; \textbf{Rhythm-Disorder Cluster} screening via Two-Stage Hybrid Screening (Rate Disorders: Bradycardia/Tachycardia via deterministic HR rules; Rhythm Arrhythmias: Normal vs. Ectopic [PVC/PAC] vs. Irregular/AFib via explainable edge ML [\texttt{LR}, \texttt{RF}, \texttt{GBDT/XGBoost}] benchmarked against \texttt{Hybrid CNN-BiLSTM-Attention}); patient-stratified clinical benchmark reruns across MIT-BIH DS1/DS2, INCART, PTB-XL, CPSC 2018, and Chapman-Shaoxing; coupled Van der Pol numerical integration (RK45) as a controlled synthetic arrhythmic SCG test-signal simulator; quantitative comparison against published public smartwatch PPG benchmarks; and structured usability testing (SUS, setup time, first-attempt capture rate).
		\item \textit{Exclusions:} Non-cardiac autonomic stress classification (e.g., WESAD affective states); 12-lead ECG morphological diagnostics such as ST-segment elevation analysis for acute myocardial infarction or structural valvular imaging, which require multi-vector electrical leads or echocardiography rather than resting chest timing; invasive cardiovascular monitoring; and continuous ambulatory exercise tracking (e.g., walking or running), where gross locomotion overwhelms micro-$g$ cardiac vibrations.
	\end{itemize}
	
	\item \textbf{What you cannot do or exclude from this project:}
	\begin{itemize}
		\item \textit{Non-diagnostic disclaimer:} The proposed prototype cannot provide clinical diagnoses, prescribe medications, or execute medical interventions; it is strictly an engineering research and screening-support prototype.
		\item \textit{Regulatory certification:} Formal FDA or CE medical device clearance is outside the scope of an undergraduate capstone timeline.
		\item \textit{Demographic boundaries:} Pediatric and neonatal populations are excluded due to substantially different baseline heart-rate ranges; evaluation is restricted to adult physiological ranges ($18\text{--}65\text{ years}$).
	\end{itemize}
	
	\item \textbf{Limitations:}
	\begin{itemize}
		\item \textit{Sensor Noise Floor:} Consumer smartphone MEMS accelerometers exhibit higher intrinsic thermal noise and quantization error than clinical-grade seismocardiometers.
		\item \textit{Posture Dependence:} Reliable sternal SCG capture requires resting supine (lying flat) or semi-recumbent posture so the accelerometer's $Z$-axis remains perpendicular to ventricular ejection recoil without hand-holding tremor.
		\item \textit{Operating System Jitter:} Mobile operating systems (iOS/Android) introduce non-deterministic thread scheduling and sampling jitter, requiring uniform cubic-spline resampling prior to spectral estimation.
		\item \textit{Smartwatch Complementarity:} Unlike wrist wearables that track continuously during ambulation, smartphone sternal SCG is an on-demand resting screening tool; optional smartwatch partnership is framed strictly as a future Complementary Dual-Mode Pathway.
	\end{itemize}
	
	\item \textbf{Ranges of data, date, and time:}
	\begin{table}[!ht]
		\footnotesize
		\centering
		\caption{Operational Ranges and System Specifications}
		\label{tab:system_ranges}
		\begin{tabular}{|p{4.2cm}|p{8.8cm}|}
			\hline
			\textbf{Parameter} & \textbf{Operational Specification} \\ \hline
			Physiological Heart Rate Range & $40\text{ bpm}$ to $180\text{ bpm}$ (Bradycardia $<60\text{ bpm}$; Normal $60\text{--}100\text{ bpm}$; Tachycardia $>100\text{ bpm}$) \\ \hline
			Peak-to-Peak Interval (PPI) Window & $333\text{ ms}$ to $1500\text{ ms}$ (physiological outlier rejection boundaries) \\ \hline
			MEMS Accelerometer Sampling Rate & $100\text{ Hz}$ to $200\text{ Hz}$ (3-axis IMU) \\ \hline
			Measurement Window Durations & $30\text{ s}$ for instantaneous pulse rate; $60\text{ s}$ to $120\text{ s}$ for statistical PRV and spectral analysis \\ \hline
			Signal-to-Noise Ratio (SNR) Threshold & $\ge 12\text{ dB}$ for valid SCG acceptance and progression to screening \\ \hline
			Edge Efficiency \& Usability Targets & Latency $< 500\text{ ms}$; Edge Model $< 2\text{ MB}$; Setup $< 60\text{ s}$; $\text{SUS} \ge 75$; Capture $\ge 90\%$ \\ \hline
			Academic Schedule & Two semesters, Academic Year 2026--2027 (see Table~\ref{tab: your-table}) \\ \hline
		\end{tabular}
	\end{table}
	
	\item \textbf{Platform(s)/Device(s) used specifically in this project:}
	\begin{itemize}
		\item \textit{Mobile Sensor Hardware:} Standard iOS (iPhone with \texttt{CoreMotion}) and Android (with \texttt{SensorManager} API) smartphones equipped with built-in 3-axis MEMS accelerometers.
		\item \textit{Mobile Frontend Application:} Cross-platform Progressive Web App (PWA) and Flutter client providing native sensor subscriptions, posture/stillness guidance, and real-time WebSocket telemetry.
		\item \textit{Computational Engine:} Python 3.9+ backend utilizing \texttt{NumPy}, \texttt{SciPy}, \texttt{Scikit-Learn}, \texttt{XGBoost}, and \texttt{PyTorch} (for the clinical CNN-BiLSTM-Attention benchmark comparator).
		\item \textit{Simulation Environment:} Python-based Runge--Kutta ODE solver (\texttt{scipy.integrate.solve\_ivp} RK45) for controlled synthetic arrhythmic SCG generation.
	\end{itemize}
\end{enumerate}

\section{Expected Results}
\begin{enumerate}
	\item \textbf{Real-Time Sternal Acquisition, Quality Gating, and Usability (Efficiency \& User-Friendliness):} The proposed mobile pipeline will acquire resting lower-sternal SCG signals with an end-to-end processing latency $< 500\text{ ms}$, achieve $100\%$ rejection of stationary inanimate-surface and active-motion artifacts via the physiological energy gate and SNR check, and achieve guided user setup in $< 60\text{ s}$ with $\ge 90\%$ first-attempt valid-signal capture and $\text{SUS} \ge 75$.
	\item \textbf{Accurate Heart-Rate Extraction and Inter-Subject Differentiation (Accuracy):} Against simultaneous reference heart-rate measurements, the DSP engine will target a Mean Absolute Error ($\text{MAE}$) $< 3.0\text{ bpm}$ and Root Mean Squared Error ($\text{RMSE}$) $< 4.5\text{ bpm}$---competitive with or outperforming published resting smartwatch PPG benchmarks---while demonstrating repeatable PRV separability between at least two distinct healthy subjects.
	\item \textbf{Two-Stage Hybrid Rhythm-Disorder Screening (Accuracy \& Edge Efficiency):} Evaluated under strict patient-stratified validation on MIT-BIH DS1/DS2 \cite{goldberger2000physiobank, dechazal2004automatic} and cross-database clinical benchmarks (INCART, PTB-XL, CPSC 2018, Chapman-Shaoxing), the deterministic rate rules ($<60\text{ bpm}$ Bradycardia, $>100\text{ bpm}$ Tachycardia) paired with explainable edge classifiers (\texttt{LR}, \texttt{RF}, \texttt{GBDT/XGBoost}, footprint $< 2\text{ MB}$) and benchmarked against \texttt{Hybrid 1D-CNN + BiLSTM + Bahdanau Attention} will target overall $\text{Accuracy} \ge 85\%$, $\text{Sensitivity} \ge 88\%$, $\text{Specificity} \ge 85\%$, and $\text{F1-score} \ge 0.86$, with feature importance rankings confirming reliance on physiologically valid interval-irregularity metrics (e.g., RMSSD, SDNN, pNN50).
	\item \textbf{Controlled Arrhythmic SCG Verification and Smartwatch Benchmark Comparison:} The coupled Van der Pol simulator (RK45) will generate reproducible synthetic arrhythmic SCG test sequences bridging real smartphone beat templates with clinical rhythm timing under MEMS noise, and a quantitative benchmark table will contextualize smartphone sternal SCG accuracy, cost, and central mechanical fidelity directly against published public smartwatch PPG data.
\end{enumerate}

\section{Timeline}
\newcolumntype{L}[1]{>{\raggedright\let\newline\\\arraybackslash\hspace{0pt}}m{#1}}
\newcolumntype{C}[1]{>{\centering\let\newline\\\arraybackslash\hspace{0pt}}m{#1}}
\newcolumntype{R}[1]{>{\raggedleft\let\newline\\\arraybackslash\hspace{0pt}}m{#1}}	

\begin{table}[!ht]
	\footnotesize
	\sloppy
	\centering
	\caption{Project Timeline}
	\label{tab: your-table}
	\begin{tabular}{|p{2.15cm}|c|c|c|c|c|c|c|c|c|c|c|c|}
		\hline
		\multicolumn{1}{|c|}{}& \multicolumn{12}{c|}{\textbf{Timeline}} \\ \cline{2-13} 
		\multicolumn{1}{|c|}{}& \multicolumn{3}{c|}{\textbf{2026}} & \multicolumn{9}{c|}{\textbf{2027}} \\ \cline{2-13} 
		\multicolumn{1}{|c|}{\multirow{-3}{2cm}{\textbf{Plan}}} & \textbf{Oct} & \textbf{Nov} & \textbf{Dec} & \textbf{Jan} & \textbf{Feb} & \textbf{Mar} & \textbf{Apr} & \textbf{May} & \textbf{Jun} & \textbf{Jul} & \textbf{Aug} & \textbf{Sep} \\ \hline
		1: Literature review \& scope & 
		\cellcolor[HTML]{000000} & 
		\cellcolor[HTML]{000000} & & & & & & & & & & \\ \hline
		2: Mobile sensing \& dataset setup & & 
		\cellcolor[HTML]{000000} & 
		\cellcolor[HTML]{000000} & & & & & & & & & \\ \hline
		3: DSP filtering \& peak detection &  & & 
		\cellcolor[HTML]{000000} & 
		\cellcolor[HTML]{000000} & & & & & & & & \\ \hline
		4: PRV feature extraction & & & & & 
		\cellcolor[HTML]{000000} & 
		\cellcolor[HTML]{000000} & & & & & & \\ \hline
		5: Hybrid rhythm-disorder model & & & & & & 
		\cellcolor[HTML]{000000} & 
		\cellcolor[HTML]{000000} & & & & & \\ \hline
		6: Controlled arrhythmia simulation & & & & & & & & 
		\cellcolor[HTML]{000000} & 
		\cellcolor[HTML]{000000} & & & \\ \hline
		7: Multi-DB, usability \& watch benchmarks & & & & & & & & & & 
		\cellcolor[HTML]{000000} & 
		\cellcolor[HTML]{000000} & \\ \hline
		8: Final thesis documentation & & & & & & & & & & & & 
		\cellcolor[HTML]{000000} \\ \hline
	\end{tabular}
\end{table}

\section{Division of Responsibilities (Two-Member Engineering Team)}
To achieve rigorous engineering execution across the full hardware-software stack, the capstone project is partitioned into two specialized engineering domains:

\begin{enumerate}
	\item \textbf{Member 1 --- Mobile Sternal SCG Acquisition, Edge DSP \& Usability Engineering:}
	\begin{itemize}
		\item \textit{Mobile Sensor Engineering \& Usability:} Will implement low-level telemetry streaming from smartphone 3-axis MEMS accelerometers at $100\text{--}200\text{ Hz}$ across iOS (\texttt{CoreMotion}) and Android (\texttt{SensorManager}), validate optimal lower-sternum placement protocols, and conduct usability trials targeting setup $< 60\text{ s}$, $\ge 90\%$ first-attempt capture, and $\text{SUS} \ge 75$.
		\item \textit{Pre-Processing \& Physiological Gating:} Will engineer the zero-phase 4th-order Butterworth bandpass filter ($0.5\text{--}40.0\text{ Hz}$), the absolute physiological energy gate ($0.005\text{ m/s}^2 \le \sigma_{raw} \le 0.50\text{ m/s}^2$), SNR verification ($\ge 12\text{ dB}$), and real-time posture and stillness guidance.
		\item \textit{Peak Detection, PRV Engine \& Stage 1--2 Validation:} Will develop the $200\text{ ms}$ moving-RMS energy envelope and adaptive sliding threshold peak detector ($\theta[n] = \mu[n] + 0.6\sigma[n]$), construct PPI tachograms, compute standardized time-domain and frequency-domain PRV metrics within a $< 500\text{ ms}$ latency budget, execute Stage 1 reference HR ($\text{MAE} < 3.0\text{ bpm}$) and Stage 2 inter-subject validation, and compile the public smartwatch PPG comparison.
	\end{itemize}

	\item \textbf{Member 2 --- Two-Stage Hybrid Rhythm-Disorder Screening \& Controlled SCG Simulation:}
	\begin{itemize}
		\item \textit{Patient-Stratified Dataset Curation \& Hybrid Architecture:} Will curate patient-stratified clinical benchmarks (MIT-BIH DS1/DS2, INCART, PTB-XL, CPSC 2018, Chapman-Shaoxing) under the Rhythm-Disorder Cluster taxonomy and implement Two-Stage Hybrid Screening: deterministic HR rules for Bradycardia/Tachycardia paired with explainable edge classifiers (Logistic Regression, Random Forest, GBDT/XGBoost; $< 2\text{ MB}$) and the \texttt{Hybrid 1D-CNN + BiLSTM + Bahdanau Attention} clinical comparator.
		\item \textit{Benchmark Rerun, Statistical Optimization \& Explainability:} Will execute patient-stratified benchmark reruns ($80\%/20\%$ and de Chazal DS1/DS2 splits, stratified 5-fold cross-validation), hyperparameter tuning, feature-importance explainability audits, and performance evaluation across Accuracy, Sensitivity, Specificity, and F1-score against the targets in Section~1.4.
		\item \textit{Controlled Arrhythmia Simulation Tool:} Will formulate and integrate the coupled nonlinear Van der Pol cardiac oscillator (\texttt{scipy.integrate.solve\_ivp} RK45) to generate controlled synthetic arrhythmic SCG test signals, bridging the scarcity of real pathological smartphone SCG recordings during end-to-end validation.
	\end{itemize}
\end{enumerate}